\chapter{Primitives and the guard}
\label{ch:primitives}

Part~II's Scheme could compute anything and touch nothing. This chapter is
about the door: the primitives that reach your computer, and the guard
that stands in front of them.

\section{The primitive table}

Every builtin you have used is a \textbf{primitive}: a Go function in a
table, with metadata. Ask for the list, or for one entry:

\begin{scm}
(length (primitives))          ;=> 111 in a bare interpreter; more in a session
(primitive-info "fs-write")
;=> ((doc . "Create a text file inside the repository or replace its entire
;            contents. Parent directories are created. For changes to an
;            existing file prefer fs-edit.")
;    (params ("path" "string" "path relative to the repository root")
;            ("content" "string" "the complete new contents"))
;    (variadic . #f) (egress . #f) (irreversible . #t) (host . #t)
;    (origin . "builtin"))
\end{scm}

The metadata is what makes the guard's job possible:

\begin{center}
\begin{tabularx}{\linewidth}{lX}
\toprule
\code{doc} & one line; shown to the model as the tool's description \\
\code{params} & \code{(name type doc)} for each argument, in order; \code{type} is one of \code{string integer number boolean list alist vector any} \\
\code{variadic} & extra arguments allowed after the declared ones \\
\code{host} & touches the operating system: files, processes, network \\
\code{irreversible} & cannot be undone by \nodwyn{}: every call confirms unless approved \\
\code{egress} & may send bytes off the machine \\
\code{origin} & \code{builtin}, or \code{scheme} for a tool you registered \\
\bottomrule
\end{tabularx}
\end{center}

Host primitives also give each parameter a \textbf{role}: \code{path},
\code{source} (a file only read), \code{dest}, \code{payload},
\code{credential}, \code{argv} (a command line), \code{method} and
\code{position}. The guard reads roles, not guesses: it knows the first
argument of \code{fs-write} is a path because the table says so, and that
\code{doc-export} reads its first argument and writes its second.

\section{The host primitives}

These exist only inside a session (not in \code{nodwyn repl}):

\begin{center}
\begin{tabularx}{\linewidth}{lX}
\toprule
\code{(fs-read path [start [count]])} & a file, or a range of lines with a \code{[lines A-B of N]} header \\
\code{(fs-write path content)} & create or replace a file (irreversible, shows a diff) \\
\code{(fs-edit path old new [all])} & replace exactly one occurrence, or all (irreversible, shows a diff) \\
\code{(fs-list path)} & a directory's names; directories end with \code{/} \\
\code{(fs-find path glob)} & file names matching a glob, at most 500 \\
\code{(index-grep path regex [glob [context]])} & search contents with a regular expression \\
\code{(git-status)}, \code{(git-diff [path [staged]])} & in a git work tree: what changed, read-only, no approval \\
\code{(shell-exec cmd args)} & run a program with a list of arguments, no shell, in the folder (irreversible) \\
\code{(shell-start cmd args)} & start a program that keeps running, in the background (irreversible) \\
\code{(shell-read id [wait])}, \code{(shell-stop id)} & its new output and state; stop it \\
\code{(shell-send id text)} & type into it (asks every time) \\
\code{(doc-export source target)} & Markdown or LaTeX to \file{.docx}, \file{.odt} or \file{.html} with pandoc (irreversible, a write of \code{target}) \\
\code{(google-doc-upload path [title])} & send a document to Google Drive as a Google Doc (asks every time; only with a stored \code{google} credential) \\
\code{(scratch-dir)} & the session's private scratch directory \\
\code{(project-notes)}, \code{(project-notes-save text)} & notes kept across sessions \\
\code{(workspace)} & the folder's kind and mounted folders \\
\code{(http-request cred url method position headers body)} & a keyed request, checked against the bindings \\
\code{(credential class)} & the handle of a stored key, or \code{\#f} \\
\bottomrule
\end{tabularx}
\end{center}

A few more are for you only and never shown to the model:
\code{credential-register} (masked input into the keychain),
\code{policy-read} and \code{policy-write} (the bindings and approval
files, with a diff confirm), \code{register-check} (next chapter) and
\code{wipe}.

\begin{tryit}
In a session, try the read side, which needs no approval:
\begin{txt}
you> :eval (index-grep "." "bas")
(("matches" (("path" . "plants.txt") ("line" . 2) ("text" . "basil")))
 ("files" . 2) ("skipped" . 0) ("truncated" . #f))
you> :eval (fs-find "." "*.txt")
(("files" "plants.txt") ("count" . 1) ("truncated" . #f))
you> :eval (fs-read "plants.txt" 2 1)
"[lines 2-2 of 2]\nbasil\n"
\end{txt}
\end{tryit}

\subsection{Programs that keep running}

\code{shell-exec} waits for a command to finish. A server, a watcher or a
program you talk to never does, so it has four primitives of its own. Here
is a web server started, used and stopped, with \code{-{}-yolo} answering:

\begin{txt}
you> :eval (display (shell-start "python3" (list "-u" "-m" "http.server" "8765")))
⚑ guard confirm shell-start: shell-start is irreversible → yes
p1: running (python3 -u -m http.server 8765)
Serving HTTP on :: port 8765 (http://[::]:8765/) ...

you> :eval (display (shell-read "p1" 2))
p1: running (python3 -u -m http.server 8765)
::ffff:127.0.0.1 - - [30/Sep/2026 13:47:10] "GET /plants.txt HTTP/1.1" 200 -

you> :eval (display (shell-stop "p1"))
p1: exited 143 (stopped by SIGTERM) (python3 -u -m http.server 8765)
\end{txt}

(In between, a \code{curl} through \code{shell-exec} fetched
\file{plants.txt}; that is the request in the log.) The rules:

\begin{itemize}
  \item \code{shell-start} is checked exactly like \code{shell-exec}: the
    same confirm, the same command-line checks, the same scrubbed
    environment and \code{\$TMPDIR}. It returns the first second of output.
  \item \code{shell-read} returns what arrived since the last read, after a
    line with the state (\code{running} or \code{exited N}); with a wait
    (up to 60 seconds) it first waits for something new. Each process keeps
    its latest 1~MiB of output, and a read says when some was dropped.
  \item \code{shell-send} types into the program. It asks every time,
    because typing into a shell is running commands.
  \item \code{shell-stop} needs no approval: it only ends what this session
    started. At most four run at once, and all of them stop when the
    session closes.
\end{itemize}

\begin{tip}
Many programs buffer their output when it is not going to a terminal, so
nothing seems to happen: Python's \code{-u} above is the cure for Python,
and \code{stdbuf -oL} for many others.
\end{tip}

\subsection{Globs}

\code{index-grep} and \code{fs-find} walk the folder themselves, in Go:
hidden entries (\file{.git}, \file{.venv}\ldots) and symbolic links are
skipped, and no ignore file is read. A glob narrows what they look at, with
the rules ripgrep uses:

\begin{center}
\begin{tabularx}{\linewidth}{lX}
\toprule
\code{*.py} & no slash: matches any path component, so \code{.py} files at any depth \\
\code{tests} & every file under a \file{tests} directory, wherever it is \\
\code{src/*.go} & a slash: matches from the start of the path, one level only \\
\code{django/**/migrations/*.py} & \code{**} spans any number of directories \\
\code{*.\{go,md\}} & either alternative \\
\code{!*.txt} & everything the rest does not match \\
\bottomrule
\end{tabularx}
\end{center}

A glob with a slash is tried from the directory you searched and from the
folder's root, so \code{(fs-find "django" "**/migrations/*.py")} and
\code{(fs-find "." "django/**/migrations/*.py")} find the same files.

\section{How the guard decides}

Every primitive application goes through one Go function, which does three
things in order: checks the argument count, asks the guard, and only then
calls the primitive. The guard runs a list of \textbf{checkers}. Each one
may say \emph{allow}, \emph{confirm} or \emph{deny}, or have no opinion,
and the strictest answer wins:

\begin{center}
\begin{tabularx}{\linewidth}{lX}
\toprule
\textbf{checker} & \textbf{what it looks at} \\
\midrule
\code{meta} & the primitive's own metadata: irreversible means confirm \\
\code{repo-pin} & every path role: inside the folder, the scratch directory or a mount \\
\code{path-denylist} & credential-looking files and directories, even inside the folder \\
\code{argv-trigger} & command lines: paths outside, network tools, dangerous flags \\
\code{protected-path} & writes to \file{.nodwyn/}, \file{.git/}, \file{.github/}, \file{scripts/}, \file{Makefile}\ldots \\
\code{scratch} & writes inside the scratch directory: no need to ask \\
\code{command-allow}, \code{write-allow} & your standing approvals \\
\code{l1} & payloads that look like credentials \\
\code{l2} & provenance: outside content heading to the network or into a command \\
\code{egress}, \code{bindings} & keyed and outgoing requests against the binding table \\
your checks & anything you registered with \code{register-check} \\
\bottomrule
\end{tabularx}
\end{center}

Two rules shape how they combine:

\begin{itemize}
  \item \textbf{A standing approval outranks a confirm, never a deny.}
    Approving \code{go test} stops the question; it does not let
    \code{go test} read \file{\textasciitilde/.ssh}.
  \item \textbf{Some confirms are security confirms}, and no approval waives
    them: a write to a protected path, content that looks like a key, a
    command line built out of tool output. They need a person, every time;
    \code{-{}-yolo} refuses them.
\end{itemize}

Here is all of that in one real session with \code{-{}-yolo}:

\begin{txt}
you> :eval (shell-exec "wc" (list "-l" "plants.txt"))
⚑ guard confirm shell-exec: shell-exec is irreversible → yes
(("status" . 0) ("stdout" . "       2 plants.txt\n") ("stderr" . ""))

you> :eval (fs-write ".nodwyn/x.scm" "1")
⚑ guard confirm fs-write: fs-write is irreversible; protected path .nodwyn/x.scm:
  it holds code that later runs with your authority; changes .nodwyn/x.scm → no
error: fs-write: denied: not confirmed: ... This needs a person to approve;
  unattended runs (--yolo) never approve it

you> :eval (fs-write "notes.txt" "sk-proj-abcdefghijklmnopqrstuvwxyz0123456789ABCD")
⚑ guard confirm fs-write: fs-write is irreversible; changes notes.txt; the written
  content contains what looks like a openai-key → no
\end{txt}

The ordinary confirm was answered yes by \code{-{}-yolo}; the two security
confirms were refused, with a reason that says exactly why.

\section{Standing approvals in files}

Commands and writes you never want to be asked about can live in two
files. Your own, for every folder:

\begin{toml}
# ~/.config/nodwyn/config.toml
[policy]
allow_commands = ["go test", "go build", "git status", "git diff", "rg"]
\end{toml}

And the folder's, tracked with it:

\begin{toml}
# <folder>/.nodwyn/allow.toml
commands = ["make lint", "make test"]
writes   = ["internal/", "docs/"]      # write prefixes that never prompt
\end{toml}

Command prefixes match whole words of the argument list: \code{"go test"}
matches \code{go test ./...} and \code{/usr/local/go/bin/go test}, but not
\code{go build}. Every use is still logged, as \emph{allow} with the reason
\code{standing approval: go test}.

\begin{tip}
Keep standing approvals read-only in spirit: tests, builds, status and
diff. Nothing that commits, pushes, deletes or reaches the network. The
point of a confirm is the one time it catches something you did not expect.
\end{tip}

\section{Provenance: where a value came from}

Every value carries \textbf{provenance tags}: not what it contains, but where
it came from. File reads are tagged \code{external:file}, command output
\code{external:shell}, network responses \code{external:http}, and what an
MCP server returns \code{external:mcp}. Tags flow
through computation:

\begin{scm}
(taint-of (fs-read "plants.txt"))                        ;=> ("external:file")
(taint-of (string-append "x" (fs-read "plants.txt")))    ;=> ("external:file")
(taint-of (string-length (fs-read "plants.txt")))        ;=> ()
(tainted? "typed by you")                                ;=> #f
\end{scm}

A result that \emph{measures} its input rather than being made of it (a
length, a comparison, a predicate) carries nothing. The guard uses the tags
to refuse sending file content to a destination no binding authorises, and
to ask before running a command line built from tool output, the classic
way a malicious file tries to get itself executed.

\code{wipe} returns the same value without its tags. It is user-only: the
model can't call it, and neither can the loop's strategy, which handles
every tool result. Provenance is not a wall against you; it is a way to
make sure that what \emph{you} did not look at does not travel by itself.

\begin{careful}
Provenance is honest about its limits: a value rebuilt one character at a
time through primitives that measure (lengths, comparisons) loses its
tags. The real lines are the binding table, fail-closed defaults and
redaction; provenance is a second net under them.
\end{careful}
